% id: 201-14
% book: 베이직쎈 대수 · 12 수학적 귀납법
% section: 유형 124 수열의 귀납적 정의의 활용
% figure: none
% answer: 
% answer_source: 
% variant_level: 0 (원본 전사)
% 이 파일은 build.mjs 가 items.json 에서 만든다. 고칠 때는 items.json 을 고친다.
\smash{\raisebox{1.5mm}{\dmkichul{베이직쎈 대수 201쪽 14번}}}
어떤 미생물은 1시간마다 2마리씩 죽고, 나머지는 각각 5마리로 분열한다. 이 미생물이 4마리였던 때로부터 $n$시간이 지난 후 살아 있는 미생물의 수를 $a_n$이라 할 때, 다음을 구하시오.
\dmsubprob{1} $a_1$
\dmsubprob{2} $a_n$과 $a_{n+1}$ 사이의 관계식
